%%%%%%%%%%%%%%%%%%%%%%%%%%%%%%%%%%%%%%%%%%%%%%%%%%%%%%%%%%%%%%%%%%%%%%%%%%%%%%%%
%%% example.tex -- how to use these TikZ figures in your own document         %%%
%%%                                                                          %%%
%%% Every <name>.tex in this folder is a bare `tikzpicture` (no preamble).   %%%
%%% To use one, your document's preamble needs three things:                 %%%
%%%                                                                          %%%
%%%   1. TikZ:              \usepackage{tikz}                                 %%%
%%%   2. the local          \usetikzlibrary{stap.components}                 %%%
%%%      tikz library       (the file tikzlibrarystap.components.code.tex    %%%
%%%                          must be on TeX's search path, e.g. in the same  %%%
%%%                          folder as your document)                        %%%
%%%   3. the colors + macros the figures expect: %%%%%%%%%%%%%%%%%%%%%%%%%%%%%%%%%%%%%%%%%%%%%%%%%%%%%%%%%%%%%%%%%%%%%%%%%%%%%%%%
%%% Colors and macros expected by the TikZ figures in this folder.           %%%
%%% Loaded by figures.tex. Replace these PLACEHOLDERS with the real          %%%
%%% definitions from your paper's preamble.                                  %%%
%%%%%%%%%%%%%%%%%%%%%%%%%%%%%%%%%%%%%%%%%%%%%%%%%%%%%%%%%%%%%%%%%%%%%%%%%%%%%%%%

%--- colors expected by the figures (fill=typeone ... typefour) ------------
\definecolor{typeone}{HTML}{7BA9DA}     % Low-degree (soft blue)
\definecolor{typetwo}{HTML}{B899C9}     % Low-degree equivalence (soft purple)
\definecolor{typethree}{HTML}{DDAA66}   % Split-and-lookup (soft orange)
\definecolor{typefour}{HTML}{7EC8A4}    % Power Residue (soft green)

%--- macros expected by the figures ----------------------------------------
\newcommand{\M}{\mathsf{M}}                     % (MDS) matrix
\newcommand{\R}{\mathsf{R}}                     % number of rounds

% Poseidon/Poseidon2/Neptune specifics
\newcommand{\Mext}{\M_{\mathcal{E}}}    % external matrix
\newcommand{\Mint}{\M_{\mathcal{I}}}    % internal matrix
\newcommand{\Rf}{\R_\mathsf{f}}                           % number of initial/final full rounds
\newcommand{\RP}{\R_\mathsf{P}}                           % number of partial rounds
\newcommand{\AddRC}{\texttt{AddRoundConstants}($\cdot$)} % round-constant layer label

% Anemoi specifics
\newcommand{\openFlystel}{\mathsf{F}_o}         % open Flystel S-box
\newcommand{\Qgamma}{\mathsf{F}_1}              % Flystel component
\newcommand{\Qdelta}{\mathsf{F}_2}              % Flystel component
\newcommand{\Mx}{\M_x}                          % Left branch matrix 
\newcommand{\My}{\M_y}                          % Right branch matrix

% RC/Monolith specifics
\newcommand{\concrete}{\texttt{Concrete}}    
\newcommand{\bricks}{\texttt{Bricks}}           
\newcommand{\bars}{\texttt{Bars}}             
\newcommand{\barslu}{{\texttt{Bar}}}  % Bar lookup element
       %%%
%%%      (copy its contents into your own preamble, or \input the file).     %%%
%%%                                                                          %%%
%%% Then drop a figure in with \input{<name>} wherever you want it.          %%%
%%%                                                                          %%%
%%% Compile this example directly to see it in action:                       %%%
%%%   pdflatex example.tex                                                   %%%
%%%                                                                          %%%
%%% NOTE: This file is only an example and is ignored by the Makefile.       %%%
%%% To build the figures themselves as individual cropped PDFs, use `make`   %%%
%%% (run `make help` for the options).                                       %%%
%%%%%%%%%%%%%%%%%%%%%%%%%%%%%%%%%%%%%%%%%%%%%%%%%%%%%%%%%%%%%%%%%%%%%%%%%%%%%%%%

\documentclass{article}

\usepackage{amsmath}
\usepackage{amssymb}

% (1) + (2): TikZ and the local component library shipped in this folder
\usepackage{tikz}
\usetikzlibrary{stap.components}

% (3): colors (typeone..typefour) and macros (\M, \openFlystel, ...) the
%      figures rely on. Merge these into your own preamble as needed.
%%%%%%%%%%%%%%%%%%%%%%%%%%%%%%%%%%%%%%%%%%%%%%%%%%%%%%%%%%%%%%%%%%%%%%%%%%%%%%%%
%%% Colors and macros expected by the TikZ figures in this folder.           %%%
%%% Loaded by figures.tex. Replace these PLACEHOLDERS with the real          %%%
%%% definitions from your paper's preamble.                                  %%%
%%%%%%%%%%%%%%%%%%%%%%%%%%%%%%%%%%%%%%%%%%%%%%%%%%%%%%%%%%%%%%%%%%%%%%%%%%%%%%%%

%--- colors expected by the figures (fill=typeone ... typefour) ------------
\definecolor{typeone}{HTML}{7BA9DA}     % Low-degree (soft blue)
\definecolor{typetwo}{HTML}{B899C9}     % Low-degree equivalence (soft purple)
\definecolor{typethree}{HTML}{DDAA66}   % Split-and-lookup (soft orange)
\definecolor{typefour}{HTML}{7EC8A4}    % Power Residue (soft green)

%--- macros expected by the figures ----------------------------------------
\newcommand{\M}{\mathsf{M}}                     % (MDS) matrix
\newcommand{\R}{\mathsf{R}}                     % number of rounds

% Poseidon/Poseidon2/Neptune specifics
\newcommand{\Mext}{\M_{\mathcal{E}}}    % external matrix
\newcommand{\Mint}{\M_{\mathcal{I}}}    % internal matrix
\newcommand{\Rf}{\R_\mathsf{f}}                           % number of initial/final full rounds
\newcommand{\RP}{\R_\mathsf{P}}                           % number of partial rounds
\newcommand{\AddRC}{\texttt{AddRoundConstants}($\cdot$)} % round-constant layer label

% Anemoi specifics
\newcommand{\openFlystel}{\mathsf{F}_o}         % open Flystel S-box
\newcommand{\Qgamma}{\mathsf{F}_1}              % Flystel component
\newcommand{\Qdelta}{\mathsf{F}_2}              % Flystel component
\newcommand{\Mx}{\M_x}                          % Left branch matrix 
\newcommand{\My}{\M_y}                          % Right branch matrix

% RC/Monolith specifics
\newcommand{\concrete}{\texttt{Concrete}}    
\newcommand{\bricks}{\texttt{Bricks}}           
\newcommand{\bars}{\texttt{Bars}}             
\newcommand{\barslu}{{\texttt{Bar}}}  % Bar lookup element


\begin{document}

% A figure placed in a float, scaled to the column width.
\begin{figure}[ht]
  \centering
  \resizebox{\linewidth}{!}{\begin{tikzpicture}[
    node distance=-0.4pt and 0.6cm, 
    sbox/.append style={fill=typeone}
]

    \node (in2) {$x_2$};

    % Invisible height reference: each column spans the 3-bar wire stack
    % (bars are centered on the in2 row at y = +0.8 / 0 / -0.8).
    \node[sbox, draw=none, fill=none] (reftop) at (0, 0.8) {};
    \node[sbox, draw=none, fill=none] (refbot) at (0,-0.8) {};

    \node[layerbox={(reftop)(refbot)}{\rotatebox{90}{\concrete}}, right=of in2] (c1) {};

    \node[layerbox={(reftop)(refbot)}{\rotatebox{90}{\bricks}}, right=of c1, fill=typeone] (b1) {};

    \node[layerbox={(reftop)(refbot)}{\rotatebox{90}{\concrete}}, right=of b1] (c2) {};

    \node[layerbox={(reftop)(refbot)}{\rotatebox{90}{\bricks}}, right=of c2, fill=typeone] (b2) {};

    \node[layerbox={(reftop)(refbot)}{\rotatebox{90}{\concrete}}, right=of b2] (c3) {};

    \node[layerbox={(reftop)(refbot)}{\rotatebox{90}{\bricks}}, right=of c3, fill=typeone] (b3) {};

    \node[layerbox={(reftop)(refbot)}{\rotatebox{90}{\concrete}}, right=of b3] (c4) {};

    \node[sbox, right=of c4, fill=typethree] (bar2) {\barslu};
    \node[sbox, above=of bar2, fill=typethree] (bar1) {\barslu};
    \node[sbox, below=of bar2, fill=typethree] (bar3) {\barslu};

    \node[layerbox={(reftop)(refbot)}{\rotatebox{90}{\concrete}}, right=of bar2] (c5) {};

    \node[layerbox={(reftop)(refbot)}{\rotatebox{90}{\bricks}}, right=of c5, fill=typeone] (b4) {};

    \node[layerbox={(reftop)(refbot)}{\rotatebox{90}{\concrete}}, right=of b4] (c6) {};

    \node[layerbox={(reftop)(refbot)}{\rotatebox{90}{\bricks}}, right=of c6, fill=typeone] (b5) {};

    \node[layerbox={(reftop)(refbot)}{\rotatebox{90}{\concrete}}, right=of b5] (c7) {};

    \node[layerbox={(reftop)(refbot)}{\rotatebox{90}{\bricks}}, right=of c7, fill=typeone] (b6) {};

    \node[layerbox={(reftop)(refbot)}{\rotatebox{90}{\concrete}}, right=of b6] (c8) {};

    \node (in1) at (in2 |- bar1) {$x_1$};
    \node (in3) at (in2 |- bar3) {$x_3$};

    \node[right=of c8] (out2) {$y_2$};
    \node (out1) at (out2 |- bar1) {$y_1$};
    \node (out3) at (out2 |- bar3) {$y_3$};

    % Wires
    \draw[wire] (in1.east |- bar1) -- (c1.west |- bar1);
    \draw[wire] (in2.east |- bar2) -- (c1.west |- bar2);
    \draw[wire] (in3.east |- bar3) -- (c1.west |- bar3);

    \draw[wire] (c1.east |- bar1) -- (b1.west |- bar1);
    \draw[wire] (c1.east |- bar2) -- (b1.west |- bar2);
    \draw[wire] (c1.east |- bar3) -- (b1.west |- bar3);

    \draw[wire] (b1.east |- bar1) -- (c2.west |- bar1);
    \draw[wire] (b1.east |- bar2) -- (c2.west |- bar2);
    \draw[wire] (b1.east |- bar3) -- (c2.west |- bar3);

    \draw[wire] (c2.east |- bar1) -- (b2.west |- bar1);
    \draw[wire] (c2.east |- bar2) -- (b2.west |- bar2);
    \draw[wire] (c2.east |- bar3) -- (b2.west |- bar3);

    \draw[wire] (b2.east |- bar1) -- (c3.west |- bar1);
    \draw[wire] (b2.east |- bar2) -- (c3.west |- bar2);
    \draw[wire] (b2.east |- bar3) -- (c3.west |- bar3);

    \draw[wire] (c3.east |- bar1) -- (b3.west |- bar1);
    \draw[wire] (c3.east |- bar2) -- (b3.west |- bar2);
    \draw[wire] (c3.east |- bar3) -- (b3.west |- bar3);

    \draw[wire] (b3.east |- bar1) -- (c4.west |- bar1);
    \draw[wire] (b3.east |- bar2) -- (c4.west |- bar2);
    \draw[wire] (b3.east |- bar3) -- (c4.west |- bar3);

    \draw[wire] (c4.east |- bar1) -- (bar1);
    \draw[wire] (c4.east |- bar2) -- (bar2);
    \draw[wire] (c4.east |- bar3) -- (bar3);

    \draw[wire] (bar1) -- (c5.west |- bar1);
    \draw[wire] (bar2) -- (c5.west |- bar2);
    \draw[wire] (bar3) -- (c5.west |- bar3);

    \draw[wire] (c5.east |- bar1) -- (b4.west |- bar1);
    \draw[wire] (c5.east |- bar2) -- (b4.west |- bar2);
    \draw[wire] (c5.east |- bar3) -- (b4.west |- bar3);

    \draw[wire] (b4.east |- bar1) -- (c6.west |- bar1);
    \draw[wire] (b4.east |- bar2) -- (c6.west |- bar2);
    \draw[wire] (b4.east |- bar3) -- (c6.west |- bar3);

    \draw[wire] (c6.east |- bar1) -- (b5.west |- bar1);
    \draw[wire] (c6.east |- bar2) -- (b5.west |- bar2);
    \draw[wire] (c6.east |- bar3) -- (b5.west |- bar3);

    \draw[wire] (b5.east |- bar1) -- (c7.west |- bar1);
    \draw[wire] (b5.east |- bar2) -- (c7.west |- bar2);
    \draw[wire] (b5.east |- bar3) -- (c7.west |- bar3);

    \draw[wire] (c7.east |- bar1) -- (b6.west |- bar1);
    \draw[wire] (c7.east |- bar2) -- (b6.west |- bar2);
    \draw[wire] (c7.east |- bar3) -- (b6.west |- bar3);

    \draw[wire] (b6.east |- bar1) -- (c8.west |- bar1);
    \draw[wire] (b6.east |- bar2) -- (c8.west |- bar2);
    \draw[wire] (b6.east |- bar3) -- (c8.west |- bar3);

    \draw[wire] (c8.east |- bar1) -- (out1);
    \draw[wire] (c8.east |- bar2) -- (out2);
    \draw[wire] (c8.east |- bar3) -- (out3);

\end{tikzpicture}
}
  \caption{The \textsc{Reinforced Concrete} permutation.}
  \label{fig:poseidon}
\end{figure}

% A figure at 60% of the text width.
\begin{figure}[ht]
  \centering
  \resizebox{0.6\linewidth}{!}{\begin{tikzpicture}[
    node distance=0.5cm and 1.5cm, 
    sbox/.append style={bbox={0.8cm}{0.8cm}},
    %state/.append style={minimum size=0.5cm},
]

    \node (x) at (0,0) {$x$};

    \node[MINUS, below=of x] (m1) {};
    \node[sbox, right=of m1, fill=typeone] (qg) {$\Qgamma$};

    \node[right=of qg] (z1) {};

    \node (y) at (z1 |- x) {$y$};

    \node[sbox, below=of qg, fill=typetwo] (e) {$x^{\frac{1}{\alpha}}$};

    \node[MINUS] (m2) at (z1 |- e) {};

    \node[sbox, below=of e, fill=typeone] (qd) {$\Qdelta$};
    \node[PLUS, left=of qd] (p) {};

    \node (z2) at (m2 |- qd) {};

    \node[below=of p] (u) {$u$};
    \node (v) at (z2 |- u) {$v$};

    \draw[-latex] (x.south) -- (m1.north);
    \draw[-latex] (y.south) -- ($(y.south |- qg)$) -- (qg.east);
    \draw[-latex] (qg.west) -- (m1.east);

    \draw[-latex] (m1.south) -- ($(m1.south |- e)$) -- (e.west);
    \draw[-latex] (m1.south) -- (p.north);

    \draw[-latex] (y.south) -- (m2.north);
    \draw[-latex] (e.east) -- (m2.west);

    \draw[-latex] (m2.south) -- ($(m2.south |- qd)$) -- (qd.east);
    \draw[-latex] (qd.west) -- (p.east);
    \draw[-latex] (m2.south) -- (v.north);

    \draw[-latex] (p.south) -- (u.north);
\end{tikzpicture}
}
  \caption{The open Flystel $\openFlystel$.}
  \label{fig:flystel-open}
\end{figure}

\end{document}
